\documentclass{article}
\usepackage[utf8]{inputenc}
\usepackage[T1]{fontenc}
\usepackage{amsmath,amssymb,amsthm}
\usepackage{hyperref}
\title{When graded PrudentBots cooperate with graded FairBots}
\author{Kodakodapa}
\date{}
\begin{document}
\maketitle
\section{Introduction}
\label{sec:graded-prudent-fair}

\subsection*{Setting}

Two programs play a one-shot prisoner's dilemma, and each can read the other's source code. A
\emph{modal agent} \cite{barasz2014robust} decides by asking what Peano arithmetic (PA) can prove about
the opponent, written with the provability box $\Box$ of the logic GL. Here $p$ means ``the
opponent cooperates with me'' and $q$ means ``the opponent cooperates with DefectBot'', the agent
that always defects. FairBot, $\Box p$, cooperates exactly when it can prove that the opponent
cooperates with it. PrudentBot adds a safety test: it also requires a proof, in the stronger
system PA+1, that the opponent would not cooperate with DefectBot, so it does not reward an
opponent that cooperates with everyone.

We study two graded families. The graded FairBot $F_k = \Box^k p$ ($k \ge 1$) asks for $k$
nested layers of provability; $F_1$ is FairBot. The graded PrudentBot
$P_j = \Box p \wedge \Box(\neg\Box^j\bot \rightarrow \neg q)$ ($j \ge 0$) runs its safety test
in PA+$j$, since $\neg\Box^{j}\bot$ expresses the consistency of PA+$(j-1)$; $P_1$ is
PrudentBot.

\subsection*{Result}

$P_j$ and $F_k$ cooperate with each other exactly when $k \le j$; when $k > j$, both defect.
In words: a PrudentBot that runs its safety test in a stronger system can afford to cooperate
with FairBots that demand more nested proof, and the exchange rate is one for one. Every $P_j$
still defects against CooperateBot, so the extra cooperation costs nothing in safety.

\subsection*{Why it matters}

Barasz et al. remark that PrudentBot's safety test may use a stronger system than PA+1 and
that this widens the set of agents it cooperates with, and online discussion of the paper
looked at specific pairs of such variants. Here the relation is pinned down exactly for the
two graded families, for all indices: the boundary of cooperation is the single inequality
$k \le j$.

\subsection*{Proved vs checked}

Proved for all $j \ge 0$ and $k \ge 1$ in Lean 4, in the height semantics of closed GL fixed
points (on the chain of heights, both agents cooperate at every height when $k \le j$; when
$k > j$ both defect above height $j + k$). The link to provability in PA is Solovay's theorem
\cite{solovay1976provability} (see also \cite{boolos1993logic}), cited, not formalised, and the reading
in the standard model assumes PA is sound.

\subsection*{Novelty statement}

This result extends \# Prior-art answer for R4 (K8: P\_j / F\_k threshold law)

Verdict: partial prior art.

\begin{itemize}
\item Barasz et al. 2014, remark after Theorem 3.2: PrudentBot's DefectBot check may use PA+n for larger n, which widens the set of agents it cooperates with without losing safety. That is the idea behind graded PrudentBots, stated informally.
\item LessWrong comments on the Barasz et al. post (orthonormal, cousin\_it) discuss indexed variants ``Prudent(m,n)Bot'' with specific pairs, e.g. Prudent(0,100)Bot cooperating with Prudent(0,1)Bot and Prudent(99,1000)Bot, and cousin\_it noting that a PrudentBot using PA and PA+1 does not cooperate with one using PA+1 and PA+2.
\item Not found: a general threshold law for P\_j against graded FairBots F\_k (cooperate iff k \textless{}= j, both defect above j+k), or a proof of it.
\item Suggested wording: the graded family was discussed informally (Barasz et al. remark; LessWrong comments); the general law and its Lean proof were not found.
\end{itemize}

\subsection*{Sources searched (owner, web search, 2026-10-17)}

\begin{itemize}
\item Barasz et al. 2014, Robust Cooperation in the Prisoner's Dilemma (arXiv 1401.5577), full text: agents CB, DB, CliqueBot, FairBot, PrudentBot, TrollBot, JustBot, WaitFairBot\_K; Theorems 3.2, 4.8, 4.10 and the obstacles section.
\item LaVictoire et al. 2014, Program Equilibrium in the Prisoner's Dilemma via Löb's Theorem (intelligence.org/files/ProgramEquilibrium.pdf).
\item LessWrong post ``Robust Cooperation in the Prisoner's Dilemma'' and its comments (orthonormal, cousin\_it, Will\_Sawin, Karl, dankane).
\item Arbital / LessWrong wiki ``Modal combat''.
\item Garrabrant, ``Prisoners' Dilemma with Costs to Modeling'' (Alignment Forum).
\item Mennen, ``Modal Combat for games other than the prisoner's dilemma'' (agentfoundations.org item 1304).
\item github.com/machine-intelligence/provability (Haskell solver and combat code; the agents folder contents were not visible).
\end{itemize}

Not searched: Critch's parametric bounded Löb paper, theses, the contents of the Haskell agents folder, non-English sources. A ``none'' verdict means ``not found in this search'', nothing more. Our prior-art search (2026-10-17) covered:

\begin{itemize}
\item Barasz et al. 2014, Robust Cooperation in the Prisoner's Dilemma (arXiv 1401.5577), full text: agents CB, DB, CliqueBot, FairBot, PrudentBot, TrollBot, JustBot, WaitFairBot\_K; Theorems 3.2, 4.8, 4.10 and the obstacles section.
\item LaVictoire et al. 2014, Program Equilibrium in the Prisoner's Dilemma via Löb's Theorem (intelligence.org/files/ProgramEquilibrium.pdf).
\item LessWrong post ``Robust Cooperation in the Prisoner's Dilemma'' and its comments (orthonormal, cousin\_it, Will\_Sawin, Karl, dankane).
\item Arbital / LessWrong wiki ``Modal combat''.
\item Garrabrant, ``Prisoners' Dilemma with Costs to Modeling'' (Alignment Forum).
\item Mennen, ``Modal Combat for games other than the prisoner's dilemma'' (agentfoundations.org item 1304).
\item github.com/machine-intelligence/provability (Haskell solver and combat code; the agents folder contents were not visible).
\end{itemize}

If you know of related work we missed, please open an issue.

\subsection*{How it was produced}

This result was found by AI research agents (the autoresearch harness). It was proved in Lean; the proof and check files are included so anyone can re-run them.

\section{Claims}

\subsection*{Claim K8}

For all j\textgreater{}=0, k\textgreater{}=1 (height-chain semantics of closed GL fixed points): P\_j=[]p\&[](\textasciitilde{}[]\textasciicircum{}jF-\textgreater{}\textasciitilde{}q) and F\_k=[]\textasciicircum{}k p cooperate at every height iff k\textless{}=j; if j\textless{}k both defect above height j+k. Lean-proved; P\_j defects vs CB.

Status: derived.

\paragraph{Verified.}
\begin{itemize}
\item \texttt{checks/K8-c046\_pjfk\_lean.py}: embedded Lean 4 proof of pj\_fk and pj\_cb compiles with no sorryAx; flipped-conclusion variant is rejected by Lean; explicit height table for j\textless{}=8,k\textless{}=8 matches (coop everywhere iff k\textless{}=j, else both defect above j+k with exact threshold) --- pass (Lean-proved, no sorryAx)
\end{itemize}

\paragraph{Evidence.}
\begin{itemize}
\item reasoning: Closed (letter-free) GL fixed points are decided by height on finite linear frames (Boolos, Logic of Provability, ch. 7), so box phi at height h means phi at all g\textless{}h; the Lean theorem is stated in exactly this semantics for arbitrary solutions.
\end{itemize}

\paragraph{Prior art.} Partial prior art.

\paragraph{Provenance.}
Run \texttt{20261017-1126}, claims K8, run-folder commit \texttt{662166ecff0651902cc5fcb04bf747c62d795af2}, published 2026-10-18.

\section{How to verify}

Each finding has its check scripts in \texttt{claims/graded-prudent-fair/checks/} of the findings repository. Run each with \texttt{python -I <script>}; it passes if it exits with status 0. The checks are Python scripts and need the packages of the research image (the image's \texttt{requirements.txt}); the Lean checks also need \texttt{lean}.

\nocite{*}
\bibliographystyle{plain}
\bibliography{refs}

\end{document}
